% Fig. (protocol): the five steps of the screening protocol as a greyscale vector flowchart, full \textwidth.
% Input from paper/sections/04_protocol.tex inside a figure* environment; needs tikz with the libraries
% arrows.meta and calc (both loaded in paper/main.tex). No fills: black text and arrows, grey frames.
% One row per step (step name left, detail right); the rows carry the detail of the protocol so that the figure
% stands alone without a numbered list in the text.
\begin{tikzpicture}[
    font=\scriptsize,
    every node/.style={outer sep=0pt},
    frame/.style={draw=black!60, line width=0.5pt},
    rule/.style={draw=black!60, line width=0.4pt},
    flow/.style={-{Stealth[length=1.5mm,width=1.2mm]}, line width=0.6pt, color=black},
  ]
  \pgfmathsetlengthmacro{\hw}{0.17\textwidth}%       heading cell
  \pgfmathsetlengthmacro{\pad}{3pt}%
  \pgfmathsetlengthmacro{\bw}{\textwidth-\hw-2*\pad-0.5pt}% body text width
  \pgfmathsetlengthmacro{\vgap}{0.30cm}%
  \pgfmathsetlengthmacro{\hhw}{0.5*\hw}%
  % \protostep{name}{north-west corner of the body cell}{heading}{body}
  \long\def\protostep#1#2#3#4{%
    \node[anchor=north west, text width=\bw, inner xsep=\pad, inner ysep=2.6pt, align=left, execute at begin node={\hyphenpenalty=10000 \exhyphenpenalty=10000 \relax}] (#1b) at #2 {#4};%
    \draw[frame] ($(#1b.north west)-(\hw,0)$) rectangle (#1b.south east);%
    \draw[rule] (#1b.north west) -- (#1b.south west);%
    \node[anchor=west, inner sep=0pt, font=\scriptsize\bfseries] at ($(#1b.west)-(\hw,0)+(\pad,0)$) {#3};%
    \coordinate (#1t) at ($(#1b.north west)-(\hhw,0)$);%
    \coordinate (#1s) at ($(#1b.south west)-(\hhw,0)$);%
  }%
  \protostep{s1}{(\hw,0)}{1\hspace{0.5em}Snapshots}%
    {A library of frozen deadlocks, each a rest state saved from a whole run: lone keystones, second traps and pair standoffs.}
  \protostep{s2}{($(s1b.south west)-(0,\vgap)$)}{2\hspace{0.5em}Probes}%
    {Each perturbation is applied to the same snapshots under each probe: quench (restart under the perturbation); replay (a copy displaced by $10^{-9}$\,m); time step $\Delta t$ changed, with per-step and with $\sqrt{\Delta t}$-scaled noise; rest-drive push $\pm$ (a constant force equal to the rest drive of force-coupled ORCA, towards ($+$) or away from ($-$) the door); lateral tilt $\pm b$; corrected implementation (the RVO2 solver).}
  \protostep{s3}{($(s2b.south west)-(0,\vgap)$)}{3\hspace{0.5em}Signatures}%
    {The responses of each perturbation to the probes, paired by snapshot: freed fraction at horizon $T$, times to free, outcome changes under replay and handedness under tilt.}
  \protostep{s4}{($(s3b.south west)-(0,\vgap)$)}{4\hspace{0.5em}Class}%
    {The class whose signatures the perturbation shows (Table~\ref{tab:signatures}): Prevent; Reshape, axial or lateral; Heat; Glitch.}
  \protostep{s5}{($(s4b.south west)-(0,\vgap)$)}{5\hspace{0.5em}Compare}%
    {The resolver (the mechanism under test) against its surrogate, a perturbation of another class whose one parameter is set so that it is fraction-matched at $T$ on the same deadlocks. The two share a mechanism only if no probe separates them (paired tests, Holm-corrected across probes).}
  \foreach \a/\b in {s1/s2, s2/s3, s3/s4, s4/s5}{\draw[flow] (\a s) -- (\b t);}%
\end{tikzpicture}%
